\documentclass[10pt, twocolumn, letterpaper]{article}

% --- Essential Packages ---
\usepackage[utf8]{inputenc}
\usepackage[T1]{fontenc}
\usepackage{libertine}
\usepackage[libertine]{newtxmath}
\usepackage{microtype}
\usepackage[hmargin=1.8cm, vmargin=2.0cm, columnsep=0.6cm]{geometry}
\usepackage{titlesec}
\usepackage{booktabs}
\usepackage{graphicx}
\usepackage{xcolor}
\usepackage{hyperref}
\usepackage{enumitem}
\usepackage{caption}
\usepackage{amsmath}

% --- Aesthetic Customization ---
\definecolor{DeepTeal}{HTML}{004D40}

\hypersetup{
    colorlinks=true,
    linkcolor=DeepTeal,
    citecolor=DeepTeal,
    urlcolor=DeepTeal
}

\titleformat{\section}
  {\color{DeepTeal}\normalfont\large\bfseries\scshape}{\thesection}{1em}{}[{\titlerule[0.5pt]}]
\titleformat{\subsection}
  {\color{DeepTeal}\normalfont\normalsize\bfseries}{\thesubsection}{1em}{}

\captionsetup{font=small, labelfont=bf, skip=4pt}
\setlength{\abovecaptionskip}{4pt}
\setlength{\belowcaptionskip}{0pt}

% --- Document Info ---
\title{
    \vspace{-1.5em}
    \rule{\linewidth}{1.5pt} \\
    \Large \textbf{CS-233: Introduction to Machine Learning} \\
    \normalsize \textit{Project Milestone 1 Report} \\
    \vspace{-0.5em}
    \rule{\linewidth}{1.5pt}
}

\author{
    \begin{tabular}{c c c}
        \textbf{Author 1} & \textbf{Author 2} & \textbf{Author 3} \\
        SCIPER: 123456 & SCIPER: 234567 & SCIPER: 345678 \\
        \textit{name1@epfl.ch} & \textit{name2@epfl.ch} & \textit{name3@epfl.ch}
    \end{tabular}
}
\date{}

\begin{document}

\maketitle

\section{Introduction}
We use the Gaming and Mental Health dataset (\texttt{features.npz}, 2000 samples, 13 features) for two tasks:
\begin{itemize}[noitemsep, nolistsep, leftmargin=*]
    \item \textbf{Regression:} predict addiction score in $[0,10]$.
    \item \textbf{Classification:} predict addiction level \{Low, Medium, High\}.
\end{itemize}
Methods implemented from scratch: Linear Regression, multiclass Logistic Regression, and KNN.

\section{Method}
\subsection{Data Preparation}
Common pipeline for all methods:
\begin{itemize}[noitemsep, nolistsep, leftmargin=*]
    \item \textbf{Split:} 70\%/30\% train-validation (seed 100). Final training uses all 1600 training samples.
    \item \textbf{Standardization:} training mean/std only; zero std replaced by 1.
    \item \textbf{Bias term:} prepend ones for Linear/Logistic Regression.
\end{itemize}

\subsection{Linear Regression}
Closed-form solution with pseudo-inverse:
\[
    \mathbf{w} = \mathbf{X}^{\dagger}\mathbf{y}, \qquad \hat{\mathbf{y}} = \mathbf{X}\mathbf{w}.
\]
\texttt{np.linalg.pinv} improves numerical stability.

\subsection{Logistic Regression}
For $C{=}3$ classes, we train softmax regression with full-batch gradient descent:
\[
    P_{i,c}=\frac{e^{Z_{i,c}-\max_c Z_{i,c}}}{\sum_{c'}e^{Z_{i,c'}-\max_{c'} Z_{i,c'}}},
\]
\[
\mathbf{W}\leftarrow\mathbf{W} - \frac{\eta}{N}\mathbf{X}^{\top}(\mathbf{P}-\mathbf{Y}).
\]
The row-max shift avoids numerical overflow.

\subsection{K-Nearest Neighbors}
KNN uses Euclidean distance:
\begin{itemize}[noitemsep, nolistsep, leftmargin=*]
    \item \textbf{Classification:} majority vote of $k$ neighbors.
    \item \textbf{Regression:} average target of $k$ neighbors.
\end{itemize}
Best $k$ is selected on validation, then retrained on full training data.

\section{Experiments and Results}
\subsection{Hyperparameter Search}
Metrics: MSE (regression), Accuracy and Macro-F1 (classification). Search grids:
\begin{itemize}[noitemsep, nolistsep, leftmargin=*]
    \item \textbf{KNN:} $k\in\{1,3,5,7,9,11,15,21\}$.
    \item \textbf{Logistic Regression:} $\eta\in\{10^{-3},3{\cdot}10^{-3},10^{-2},3{\cdot}10^{-2},10^{-1},3{\cdot}10^{-1}\}$, iterations $\in\{200,500,1000,2000\}$.
\end{itemize}

Figures~\ref{fig:knn} and \ref{fig:logreg} summarize validation behavior and selected hyperparameters.

\begin{figure}[t]
\centering
\includegraphics[width=\linewidth]{figures/knn_tuning.png}
\caption{KNN validation curves. Best overall point is $k{=}9$ for both classification (Acc/Macro-F1) and regression (lowest MSE).}
\label{fig:knn}
\end{figure}

\begin{figure}[t]
\centering
\includegraphics[width=\linewidth]{figures/logreg_heatmap.png}
\caption{Logistic Regression grid search. $\eta{=}0.3$ with 500--1000 iterations gives the best validation Acc and Macro-F1.}
\label{fig:logreg}
\end{figure}

\begin{figure}[t]
\centering
\includegraphics[width=0.88\linewidth]{figures/confusion_matrix.png}
\caption{Confusion matrix of best Logistic Regression model on test set ($\eta{=}0.3$, 1000 iterations). Main confusion appears between Medium and High.}
\label{fig:cm}
\end{figure}

\begin{figure}[t]
\centering
\includegraphics[width=\linewidth]{figures/class_distribution.png}
\caption{Class distribution for train/test splits. Low class is dominant; Macro-F1 is therefore reported with Accuracy.}
\label{fig:dist}
\end{figure}

\subsection{Final Comparison}

\begin{table}[h]
\centering
\caption{Final method comparison on validation and held-out test sets.}
\label{tab:results}
\setlength{\tabcolsep}{4pt}
\begin{tabular}{@{}llcc@{}}
\toprule
\textbf{Method} & \textbf{Metric} & \textbf{Val.} & \textbf{Test} \\
\midrule
Linear Regression      & MSE         & 0.931 & \textbf{0.994} \\
KNN Reg.\ ($k{=}9$)   & MSE         & 1.947 & 1.865 \\
\addlinespace
LogReg ($\eta{=}0.3$, 1000) & Acc.\,(\%) & \textbf{85.83} & \textbf{88.25} \\
KNN Cls.\ ($k{=}9$)   & Acc.\,(\%) & 80.63 & 79.75 \\
\addlinespace
LogReg ($\eta{=}0.3$, 1000) & Macro-F1 & 0.725 & \textbf{0.819} \\
KNN Cls.\ ($k{=}9$)   & Macro-F1 & 0.579 & 0.590 \\
\bottomrule
\end{tabular}
\end{table}

\section{Discussion and Conclusion}
\paragraph{Key outcomes.}
\begin{itemize}[noitemsep, nolistsep, leftmargin=*]
    \item \textbf{Best regressor:} Linear Regression (test MSE \textbf{0.994}).
    \item \textbf{Best classifier:} Logistic Regression (test Acc \textbf{88.25\%}, Macro-F1 \textbf{0.819}).
    \item KNN is competitive but lower on both tasks.
\end{itemize}

\paragraph{Conclusion.}
For this milestone, Linear Regression and Logistic Regression are the strongest choices. Next improvements: class-weighted loss, regularization, and feature engineering.

\end{document}
